% Width: 112 mm (text width, figure); height: 64 mm.
% Use at native size with \input; all labels are at least footnotesize.
% Population schematic, not random samples or an ICA optimization result.
% Independent sources s_1,s_2 ~ Uniform[-sqrt(3),sqrt(3)] have mean 0, Var=1.
% R(theta)=[cos(theta),-sin(theta);sin(theta),cos(theta)].
% x=R(30 deg)s, Cov(x)=I_2; z=R(-30 deg)x=s.
% The square outline is the support of the continuous population.
% Dots are the deterministic Cartesian grid {(j/2,k/2):j,k=-3,...,3}.
% Its 49 equally weighted points also have mean 0 and covariance I_2
% with denominator n, but do not constitute a random sample.
% Plot scale: 11 mm per dimensionless coordinate unit in both panels.
% Matrix convention: S,X,Z in R^{49 x 2} store samples as rows.
% X=S R(30 deg)^T, W=R(30 deg), Z=X W=S.
% Dashed arrows are columns of W; PCA has no preferred direction because
% its population eigenvalues are equal, even though x_1,x_2 are dependent.
% ICA only identifies these non-Gaussian sources up to scale and permutation.
\begin{tikzpicture}[
  x=1mm,y=1mm,font=\figurefont\footnotesize,text=FigureInk,
  ica axis/.style={-{Stealth[length=1.7mm]},draw=FigureMuted,line width=.7pt},
  ica direction/.style={-{Stealth[length=1.8mm]},draw=FigureModel,
    line width=.9pt,dashed}
]
  \path[use as bounding box] (0,0) rectangle (112,64);
  \node at (27,60) {（a）白化后：不相关};
  \node at (85,60) {（b）源坐标：独立};
  \node at (27,53) {$\bm{x}=\bm{R}(30^\circ)\bm{s}$};
  \node at (85,53) {$\bm{z}=\bm{R}(-30^\circ)\bm{x}$};

  \begin{scope}[shift={(27,28)},x=7mm,y=7mm]
    \begin{scope}[rotate=30]
      \draw[draw=FigureInput,fill=FigureInput!8!FigurePaper,line width=.8pt]
        ({-sqrt(3)},{-sqrt(3)}) rectangle ({sqrt(3)},{sqrt(3)});
      \foreach \j in {-3,...,3}{
        \foreach \k in {-3,...,3}{
          \fill[FigureInput] ({\j/2},{\k/2}) circle (1pt);
        }
      }
      \draw[ica direction] (-2.1,0)--(2.25,0)
        node[above right] {$\bm w_1$};
      \draw[ica direction] (0,-2.1)--(0,2.25)
        node[above left] {$\bm w_2$};
    \end{scope}
    \draw[ica axis] (-2.5,0)--(2.5,0) node[right] {$x_1$};
    \draw[ica axis] (0,-2.5)--(0,2.5) node[above] {$x_2$};
  \end{scope}

  \draw[book arrow,draw=FigureModel] (48,28)--(64,28)
    node[midway,above,text=FigureInk] {旋转解混};

  \begin{scope}[shift={(85,28)},x=7mm,y=7mm]
    \draw[draw=FigureOutput,fill=FigureOutput!8!FigurePaper,line width=.8pt]
      ({-sqrt(3)},{-sqrt(3)}) rectangle ({sqrt(3)},{sqrt(3)});
    \foreach \j in {-3,...,3}{
      \foreach \k in {-3,...,3}{
        \fill[FigureOutput] ({\j/2},{\k/2}) circle (1pt);
      }
    }
    \draw[ica axis] (-2.5,0)--(2.5,0) node[right] {$z_1=s_1$};
    \draw[ica axis] (0,-2.5)--(0,2.5) node[above] {$z_2=s_2$};
  \end{scope}
\end{tikzpicture}%
